% ============================================================
% CHAPTER 5 — RESULTS
% ============================================================
\chapter{Results}
\label{ch:results}

This chapter reports the empirical results of the calibration--holdout evaluation and tests the three hypotheses of Section~\ref{sec:hypotheses}. Section~\ref{sec:results-diagnostics} establishes that the MCMC inference is trustworthy; Sections~\ref{sec:results-h1}--\ref{sec:results-h3} address H1, H2, and H3 in turn; and Section~\ref{sec:results-summary} collates the verdicts.

\paragraph{Note on reproduction.}
The tables and figures in this chapter are produced by \texttt{src/\allowbreak run\_all\_models.py} and written to \texttt{outputs/\allowbreak results/} and \texttt{outputs/\allowbreak figures/}; all numeric values quoted in the prose are read directly from those generated tables. The standard BG/NBD and Gamma-Gamma models are estimated with the \texttt{pymc-marketing} implementation of the numerically stable Fader--Hardie likelihood, while the hierarchical model is a custom PyMC implementation of the same likelihood with partial pooling across country segments.

\section{MCMC Convergence and Diagnostics}
\label{sec:results-diagnostics}

Before interpreting any posterior, the quality of the sampling is verified using the diagnostics of Section~\ref{sec:mcmc}. Figure~\ref{fig:rhat} shows the split-$\widehat{R}$ values for the pooled BG/NBD, hierarchical BG/NBD, and Gamma-Gamma models. Convergence is indicated when all parameters satisfy $\widehat{R} < 1.01$; the effective sample sizes and the number of divergent transitions for each model are summarised below. The trace and pair plots (Figures~\ref{fig:trace} and~\ref{fig:pairs}) provide a visual check on mixing and posterior correlation structure.

\condfig{../outputs/figures/rhat_bgnbd_standard.png}{0.7\linewidth}{Split-$\widehat{R}$ convergence diagnostics for the pooled BG/NBD model. The dashed line marks the $1.01$ threshold.}{fig:rhat}

\condfig{../outputs/figures/trace_bgnbd_standard.png}{0.85\linewidth}{Posterior trace plots for the pooled BG/NBD population parameters $(r, \alpha, a, b)$. Well-mixed, stationary chains indicate reliable sampling.}{fig:trace}

\condfig{../outputs/figures/pairs_bgnbd.png}{0.75\linewidth}{Pairwise posterior marginals for the BG/NBD parameters, with divergences highlighted. Strong correlations would motivate reparameterisation.}{fig:pairs}

For the hierarchical model, a non-centred parameterisation (Section~\ref{sec:noncentered}) with a tight between-segment scale was adopted to control the funnel geometry that the small country segments induce; Figure~\ref{fig:rhat-hier} confirms that this resolved the sampling pathologies. Across all three models the maximum split-$\widehat{R}$ is $1.002$ and the minimum bulk effective sample size exceeds $1{,}100$, with no divergent transitions. The diagnostics therefore indicate that the reported posteriors are well converged and provide a sound basis for the inferences that follow.

\condfig{../outputs/figures/rhat_bgnbd_hierarchical.png}{0.7\linewidth}{Split-$\widehat{R}$ for the hierarchical BG/NBD model under the non-centred parameterisation.}{fig:rhat-hier}

\section{H1: Predictive Accuracy and Calibrated Uncertainty}
\label{sec:results-h1}

\subsection{Point and Ranking Accuracy}
\label{sec:results-accuracy}

Table~\ref{tab:metrics_comparison} reports the headline comparison across all models for both holdout transaction prediction and CLV, covering error (MAE, RMSE, MAPE), ranking (Spearman, Gini), and top-$k$ relevance (NDCG). H1 predicts that the Bayesian BG/NBD~$+$~Gamma-Gamma model is competitive with or superior to the RFM-heuristic and XGBoost baselines. The results show that the Bayesian model attains a transaction MAE of $1.74$ against $2.12$ for XGBoost and $2.33$ for the RFM heuristic, and a CLV Gini of $0.86$ (versus $0.80$ for XGBoost). The ranking metrics are the most decisive: the Bayesian model reaches an NDCG@100 of $0.91$ against $0.53$ for XGBoost, so it identifies the highest-value customers far more reliably than either baseline while also carrying the lowest CLV MAE (\pounds$847$ against \pounds$1{,}478$ for XGBoost). The size of this ranking gap warrants explanation, since a value near $0.9$ against one near $0.5$ is a large margin. Two factors combine to produce it. First, the generative model borrows strength across the customer base, so even customers with sparse histories receive a sensibly ordered value estimate, whereas the discriminative XGBoost, trained to minimise average error, is comparatively indifferent to the exact ordering within the small high-value tail that NDCG@100 rewards. Second, the leakage-safe design gives XGBoost a shorter effective training window (Section~\ref{sec:limitations-eval}), which handicaps precisely the top-of-ranking task; its accuracy on bulk error metrics (transaction MAE $2.12$) is much closer to the Bayesian model than the NDCG gap alone would suggest. The Bayesian model is thus superior, not merely competitive, on out-of-sample accuracy.

\resulttable{metrics_comparison}{Model comparison across transaction and CLV prediction metrics.}{tab:metrics_comparison}

The decile calibration plot (Figure~\ref{fig:calibration}) shows predicted versus actual transactions for each model. Points close to the diagonal indicate good calibration of the point predictions; systematic departure in the top deciles would signal over- or under-prediction for the most active customers.

\condfig{../outputs/figures/calibration_comparison.png}{0.95\linewidth}{Decile calibration: mean predicted versus mean actual holdout transactions for each model. The diagonal denotes perfect calibration.}{fig:calibration}

\subsection{The XGBoost Benchmark}
\label{sec:results-xgboost}

Because XGBoost is the strongest baseline, understanding \emph{why} it trails the Bayesian model is instructive. Figure~\ref{fig:xgb-importance} reports the gain-based feature importances of its transaction model. The ranking is dominated by \textbf{frequency} ($0.29$) and \textbf{purchase rate} ($0.16$), which together account for roughly $45\%$ of the total gain: given twenty-four engineered features, the model spends most of its explanatory power rediscovering the recency--frequency signal that the BG/NBD encodes structurally in its generative story. The remaining importance is spread thinly across the engineered rhythm features the parametric model cannot use, namely day-of-week purchase fractions and the mean and variability of inter-purchase times. That these behavioural extras add so little confirms that, on this dataset, the predictive content is concentrated in exactly the summaries the generative model already exploits.

\condfig{../outputs/figures/xgboost_feature_importance.png}{0.85\linewidth}{Gain-based feature importance of the XGBoost transaction model (top 12 of 24 features). Frequency and purchase rate dominate; the engineered day-of-week and inter-purchase-time features contribute comparatively little.}{fig:xgb-importance}

This explains the pattern in Table~\ref{tab:metrics_comparison}. XGBoost is genuinely competitive on bulk error: a transaction MAE of $2.12$ is far better than the naive and RFM baselines and within striking distance of the Bayesian model's $1.74$. Where it falls away sharply is at the top of the ranking, the regime that matters for targeting: its NDCG@100 is $0.53$ against $0.91$, and its CLV MAE is \pounds\num{1478} against \pounds\num{847}. Three factors combine to produce this gap. First, the squared-error objective optimises average accuracy, not the ordering of the sparse high-value tail that NDCG rewards. Second, the leakage-safe design trains it on a $21.4$-week inner window (Section~\ref{sec:xgboost}), a shorter supervision horizon than the full calibration window the Bayesian models use. Third, and most fundamentally, it has no mechanism to borrow strength across data-sparse customers and no native uncertainty, so it cannot express how confident any individual prediction is. XGBoost is thus the right benchmark precisely because it is a strong point predictor that nonetheless loses on the two axes, top-of-ranking accuracy and calibrated uncertainty, that motivate the Bayesian approach.

\subsection{Uncertainty Calibration}
\label{sec:results-uncertainty}

The distinctive claim of H1 is that the Bayesian model adds \emph{calibrated uncertainty} that point-prediction baselines cannot provide. Table~\ref{tab:coverage} reports the empirical coverage of posterior predictive credible intervals, the mean interval width, and the CRPS. Coverage close to the nominal level (e.g.\ empirical coverage near $90\%$ for the $90\%$ interval) indicates well-calibrated uncertainty; the observed coverage is $89.3\%$, almost exactly the nominal $90\%$, so the posterior predictive intervals are well calibrated. The posterior predictive distributions for a sample of customers (Figure~\ref{fig:postpred}) illustrate how the model expresses heterogeneous uncertainty across the customer base.

\resulttable{credible_interval_coverage}{Posterior predictive credible-interval coverage and CRPS.}{tab:coverage}

\condfig{../outputs/figures/posterior_predictive_bgnbd.png}{0.95\linewidth}{Posterior predictive distributions of holdout transactions for a sample of customers (violins), with realised values overlaid. The spread of each violin is the model's per-customer uncertainty.}{fig:postpred}

\subsection{Activity Prediction}
\label{sec:results-palive}

The $P(\text{alive})$ estimates are evaluated as a classifier of holdout activity (Table~\ref{tab:p_alive}). Because a customer with no repeat purchase is trivially assigned $P(\text{alive}) = 1$ under the BG/NBD model (they cannot have dropped out after a repeat that never occurred), the discrimination metrics are computed on repeat purchasers, for whom the estimate is informative. On this population the BG/NBD model achieves an AUC of $0.71$ and a Brier score of $0.21$, indicating moderate but useful discrimination of customers who remain active. An AUC of $0.71$ should be read as a genuine signal rather than a disappointing one: activity is a latent state that no model can observe, and a customer who is truly still active may simply not purchase again within the finite holdout window through ordinary randomness, which caps how well \emph{any} classifier can separate the two groups on this task. The value is well above the chance level of $0.5$, and, as the calibration-first framing of this thesis stresses, the probability is more useful for being calibrated than for being sharp. Figure~\ref{fig:palive} shows that $P(\text{alive})$ increases with calibration frequency, as theory predicts.

\resulttable{p_alive_evaluation}{$P(\text{alive})$ evaluation: binary activity-prediction metrics.}{tab:p_alive}

\condfig{../outputs/figures/p_alive_bgnbd_standard.png}{0.95\linewidth}{Distribution of $P(\text{alive})$ overall and by frequency group. Higher-frequency customers concentrate at higher $P(\text{alive})$.}{fig:palive}

\subsection{A Worked Example: One Customer End to End}
\label{sec:worked-example}

The aggregate metrics are best understood by watching the model reason about a single customer. Customer~14817, introduced in Section~\ref{sec:dist-data}, enters with only four numbers: frequency $x = 4$, recency $t_x = 33.1$ weeks, observation window $T = 49.2$ weeks, and mean repeat spend $m = \pounds 526$. From these the fitted model produces, in sequence, every quantity this thesis is concerned with.

First, activity. Despite the fifteen-week silence before the calibration cut-off, the pattern of four repeat purchases spread across the window yields $P(\text{alive}) = 0.88$: probably, but not certainly, still active. Second, future purchasing. The conditional expectation over the 40-week holdout horizon is $E[X] = 2.81$ transactions (90\% credible interval $[2.74, 2.88]$). Third, spend. The Gamma-Gamma model shrinks the customer's observed \pounds 526 toward the population mean, giving an expected value of roughly \pounds 485 per transaction. Multiplying the two components gives an expected CLV of \pounds\num{1364}, with a strikingly tight point-posterior interval of $[\pounds\num{1326}, \pounds\num{1401}]$.

That tight interval is where the worked example becomes instructive, because it reflects only \emph{parameter} uncertainty and is therefore the wrong interval for a decision about a realised outcome. The posterior \emph{predictive} distribution of the customer's actual holdout value, which additionally accounts for the sampling variability of the transaction count and the possibility of dropout, is far wider: a 90\% interval of $[\pounds 0, \pounds\num{2888})$, with substantial mass at zero (Figure~\ref{fig:worked}). The two distributions disagree sharply about the very question a targeting rule asks. At the headline intervention cost of \pounds 600, the probability $P(\text{CLV} > \pounds 600)$ computed from the posterior of the expectation is a falsely emphatic $1.00$, because the narrow spike sits entirely above the threshold; computed from the predictive distribution it is a candid $0.77$, acknowledging the roughly one-in-four chance the customer returns too little to justify the spend. This single customer previews the population-wide artefact examined in Section~\ref{sec:results-h3}. In the event, customer~14817 returned twice in the holdout and spent \pounds\num{1010}: profitable to target, above the cost, and comfortably inside the predictive interval that the point estimate's tight band would have hidden.

\condfig{../outputs/figures/worked_example_customer.png}{0.95\linewidth}{Two posteriors for one customer (ID~14817). The narrow posterior of $E[\text{CLV}]$ (orange) reflects parameter uncertainty only and lies entirely above the \pounds 600 cost, giving $P(\text{CLV}>\pounds 600)=1.00$. The posterior predictive of realised CLV (blue) is far wider, placing $\sim\!23\%$ of its mass below the cost, so the same probability computed correctly is $0.77$. The realised holdout value (\pounds\num{1010}) falls well inside the predictive distribution.}{fig:worked}

\paragraph{Verdict on H1.}
Taken together, the accuracy metrics (Table~\ref{tab:metrics_comparison}) and the uncertainty diagnostics (Table~\ref{tab:coverage}) indicate that H1 is supported: the Bayesian model is superior on accuracy, most clearly on the ranking metrics that matter for targeting, while uniquely providing well-calibrated predictive intervals.

\section{H2: Hierarchical Partial Pooling Across Segments}
\label{sec:results-h2}

H2 predicts that hierarchical partial pooling reduces prediction error for small country segments relative to the pooled (complete-pooling) model. The natural benchmark for this claim is the three-way comparison of \emph{complete} pooling, \emph{partial} pooling, and \emph{no} pooling, the last being an independent BG/NBD fitted separately on each segment. Table~\ref{tab:country_mae} reports per-country transaction MAE for all three specifications alongside the classical baselines, and Figure~\ref{fig:country-mae} visualises the comparison.

Two contrasts matter. Against complete pooling, the hierarchical model does not reduce error in the small markets: for Germany ($n=74$) its MAE is $1.635$ versus $1.622$, and for France ($n=54$) $2.318$ versus $2.223$, in both cases a marginal increase. Against no pooling, however, the pattern is exactly the one partial pooling is designed to produce. No pooling is the \emph{worst} of the three in every small segment, where the sparse per-segment data over-fit (France $2.444$, Germany $1.670$), whereas complete pooling is instead worst in the larger, behaviourally distinct ``Other'' segment ($1.695$), where imposing the global parameters ignores a genuine difference. Partial pooling is never the worst in any segment and attains the lowest error aggregated over the three non-UK segments ($1.746$, against $1.755$ for no pooling and $1.757$ for complete pooling). For the data-rich UK segment all three coincide ($\approx 1.74$), as expected where data are abundant. The differences among the specifications are small in absolute terms, a consequence of the tight between-segment prior needed for stable sampling, which shrinks the segment parameters strongly toward the global values.

\resulttable{country_level_mae}{Per-country transaction-prediction MAE by model (RQ2/H2).}{tab:country_mae}

\condfig{../outputs/figures/country_mae_comparison.png}{0.95\linewidth}{Per-country transaction MAE by model. Differences concentrated in the small segments (Germany, France, Other) test the partial-pooling hypothesis.}{fig:country-mae}

The mechanism behind any improvement is visible in the shrinkage of the segment-level parameters toward the global mean. Figure~\ref{fig:shrinkage-r} shows the per-segment posteriors for the BG/NBD shape parameter $r$: small segments exhibit wider posteriors pulled toward the global estimate, whereas the UK segment is estimated precisely and close to its no-pooling value.

\condfig{../outputs/figures/hierarchical_shrinkage_r.png}{0.7\linewidth}{Per-segment posterior of the BG/NBD parameter $r$ (90\% HDI), with the global mean marked. Small segments are shrunk toward the global estimate.}{fig:shrinkage-r}

\paragraph{Verdict on H2.}
H2 is partially supported. Partial pooling does not beat complete pooling on this relatively homogeneous country segmentation, but the borrowing-of-strength mechanism is demonstrably operating: it strictly dominates no pooling in every small segment, is never the worst specification anywhere, and attains the lowest error aggregated across the non-UK segments. It thus delivers the robustness that motivates the hierarchical model, adapting between the two pooling extremes. Two caveats qualify this reading: the between-specification differences are small in absolute terms, and the tight between-segment prior required for stable sampling itself limits how far partial pooling can depart from complete pooling (Section~\ref{sec:limitations}).

\section{H3: Decision-Theoretic Targeting}
\label{sec:results-h3}

H3 predicts that targeting customers by the posterior probability $P(\text{CLV} > c)$ yields higher cumulative realised value than targeting by point-estimate CLV. Table~\ref{tab:targeting_sim_bg_nbd_bayesian} reports the targeting simulation for the Bayesian model: cumulative net value under the point-estimate rule, the posterior-probability rule, and the oracle, across targeting depths and the swept intervention cost. Figure~\ref{fig:targeting-sim} visualises the strategy curves and the incremental advantage of the posterior rule.

\resulttable{targeting_simulation_bg_nbd_bayesian}{Decision-theoretic targeting (BG/NBD): posterior-probability versus point-estimate versus oracle net value, swept over intervention cost.}{tab:targeting_sim_bg_nbd_bayesian}

\condfig{../outputs/figures/targeting_simulation_bg_nbd_bayesian.png}{0.95\linewidth}{Decision-theoretic targeting for the Bayesian model at the headline cost of \pounds 600. Left: cumulative net value by strategy and depth. Right: the net-value difference between the posterior-probability rule and the point estimate.}{fig:targeting-sim}

The probability score is computed from the posterior \emph{predictive} distribution of realised CLV rather than from the posterior of \emph{expected} CLV. This distinction, the same one that governs interval coverage in Section~\ref{sec:results-uncertainty}, is decisive here: the posterior of the expectation reflects only parameter uncertainty and, on a dataset of this size, places $P(\text{CLV} > c)$ at essentially $0$ or $1$ for almost every customer, so the ranking degenerates into arbitrary tie-breaking and the rule spuriously appears to fail by more than $50\%$. Using the predictive distribution removes that artefact. Figure~\ref{fig:prob-dist} shows the effect across the whole customer base: computed from the posterior of the expectation, $96\%$ of customers are pinned at a probability of $0$ or $1$, leaving almost nothing for the ranking to sort on; computed from the posterior predictive, the same probability is graded across the full $[0,1]$ range, and the worked example of Section~\ref{sec:worked-example} (Figure~\ref{fig:worked}) is one customer drawn from this second, informative distribution.

\condfig{../outputs/figures/prob_expectation_vs_predictive.png}{0.98\linewidth}{Distribution of the per-customer targeting probability $P(\text{CLV}>\pounds 600)$ across all \num{4{,}522} customers. \emph{(a)} Computed from the posterior of $E[\text{CLV}]$, the probability saturates: $96\%$ of customers sit at $0$ or $1$. \emph{(b)} Computed from the posterior predictive of realised CLV, it is graded and usable as a decision score. This is the saturation artefact that makes the expectation-based rule fail for mechanical reasons.}{fig:prob-dist} Table~\ref{tab:h3_risk} makes the contrast concrete on the risk-sensitive metrics that the probability rule is meant to serve: at the headline cost and depth, the predictive probability rule matches the point estimate on both the hit rate (the share of targeted customers whose realised value exceeds the cost) and wasted spend (the cumulative shortfall on unprofitable targets), whereas the expected-CLV version of the same rule collapses on both.

\begin{table}[htbp]
  \centering
  \caption{Risk-sensitive targeting metrics at the headline cost of \pounds 600 and a targeting depth of $20\%$. \emph{Hit rate} is the fraction of targeted customers whose realised holdout value exceeds the intervention cost; \emph{wasted spend} is the cumulative shortfall on targeted customers who fall below it. The point-estimate rule and the posterior-\emph{predictive} probability rule are close on both; the probability rule computed from the posterior of \emph{expected} CLV instead collapses, illustrating the saturation artefact.}
  \label{tab:h3_risk}
  \begin{tabular}{lrr}
    \toprule
    Targeting rule & Hit rate & Wasted spend (\pounds) \\
    \midrule
    Point estimate, $\mathrm{E}[\text{CLV}]$            & $0.848$ & $55{,}040$ \\
    $P(\text{CLV} > c)$, expected-CLV posterior         & $0.619$ & $144{,}738$ \\
    $P(\text{CLV} > c)$, posterior predictive           & $0.827$ & $60{,}452$ \\
    \bottomrule
  \end{tabular}
\end{table}

At the headline intervention cost of \pounds 600, the posterior-probability rule then captures \pounds 3.82M in cumulative net value at a targeting depth of $20\%$, against \pounds 3.89M for the point-estimate rule, a shortfall of only $1.7\%$. Across the swept cost grid (\pounds 100 to \pounds 2{,}000) and across targeting depths the two rules are, in practice, equivalent: the gap is within a few percentage points everywhere and turns marginally in the probability rule's favour at the highest costs and depths, where mis-targeting a customer whose realised value falls below the intervention cost becomes costly enough for its risk-aversion to matter. This near-parity is the expected outcome on reflection: ranking by posterior mean CLV is close to optimal for maximising \emph{total} realised value, so the additional tail information the probability rule exploits is not needed for that objective. For completeness, Figure~\ref{fig:lift} reports the conventional targeting-lift curves, which measure revenue captured as a function of targeting depth.

\condfig{../outputs/figures/targeting_lift_curves.png}{0.95\linewidth}{Cumulative-gain and lift curves for all models. Steeper early rise indicates better identification of high-value customers.}{fig:lift}

\paragraph{Verdict on H3.}
The simulation indicates that H3 is not supported: targeting by $P(\text{CLV} > c)$ does not outperform point-estimate targeting for total realised value; the two are effectively equivalent. The finding is informative rather than merely negative, on two counts. First, it shows that the decision value of the full posterior depends on the objective: for maximising expected captured value the posterior mean is sufficient, whereas the probability rule expresses a risk-averse preference that would be advantageous under a loss-sensitive objective or a hard budget constraint (Section~\ref{sec:improvement}), neither of which the total-value criterion rewards. Second, it exposes a concrete pitfall in the decision-theoretic use of Bayesian models: a rule defined on the posterior of an expectation rather than on the posterior predictive can fail for purely mechanical reasons, and respecting that distinction is essential to using the posterior correctly.

\section{Summary of Hypothesis Tests}
\label{sec:results-summary}

Table~\ref{tab:verdicts} collates the verdicts. The detailed interpretation and its implications for theory and practice are discussed in Chapter~\ref{ch:conclusion}.

\begin{table}[htbp]
  \centering
  \caption{Summary of hypothesis-test outcomes.}
  \label{tab:verdicts}
  \begin{tabularx}{\linewidth}{@{}p{0.8cm} Y l@{}}
    \toprule
    \textbf{ID} & \textbf{Hypothesis (abridged)} & \textbf{Verdict} \\
    \midrule
    H1 & Bayesian model competitive/superior on accuracy with calibrated uncertainty & Supported \\
    \addlinespace
    H2 & Partial pooling reduces error for small country segments & Partially supported \\
    \addlinespace
    H3 & $P(\text{CLV} > c)$ targeting beats point-estimate targeting & Not supported \\
    \bottomrule
  \end{tabularx}
\end{table}
